\documentclass[11pt,oneside]{article}
\pdfminorversion=7
\usepackage{QuizStyle}
\hypersetup{pdftitle={20261110 Quiz for MATH230 Probability \& Statistics},pdfauthor={Jungwoo Kim},pdfsubject={Quiz for MATH230 Probability \& Statistics},pdfcreator={LaTeX}}
\begin{document}
\quiztitle{20261110 Quiz}

\begin{exercise}{7.16}
Let $X$ have a gamma distribution with shape $\alpha>0$ and scale $\theta>0$, with density
\[
f_X(x)=\frac{x^{\alpha-1}e^{-x/\theta}}{\Gamma(\alpha)\theta^\alpha}, \qquad x>0.
\]
Show that the $r$th moment about the origin is $\mu'_r=\theta^r\Gamma(\alpha+r)/\Gamma(\alpha)$ for any nonnegative integer $r$. \textit{Hint.} Substitute $u=x/\theta$.
\end{exercise}
\begin{solution}
From the definition of a moment and the given density,
\[
\begin{aligned}
\mu'_r=\Exp[X^r]
&=\frac{1}{\Gamma(\alpha)\theta^\alpha}\int_0^\infty x^{\alpha+r-1}e^{-x/\theta}\,\dd x\\
&=\frac{\theta^r}{\Gamma(\alpha)}\int_0^\infty u^{\alpha+r-1}e^{-u}\,\dd u
=\boxed{\frac{\theta^r\Gamma(\alpha+r)}{\Gamma(\alpha)}}.
\end{aligned}
\]
Here $x=\theta u$ and $\dd x=\theta\,\dd u$; the last integral is the gamma function.
\end{solution}

\begin{exercise}{7.17}
The random variable $X$ has the discrete uniform distribution
\[
\Prob(X=x)=\frac1k, \qquad x=1, 2, \ldots, k,
\]
where $k\in\N$. Show that its moment-generating function is
\[
M_X(t)=\frac{e^t(1-e^{kt})}{k(1-e^t)}\quad(t\ne0).
\]
\end{exercise}
\begin{solution}
The finite geometric-series formula gives, for $t\ne0$,
\[
\begin{aligned}
M_X(t)=\Exp[e^{tX}]
&=\frac1k\sum_{x=1}^k e^{tx}
=\frac{e^t}{k}\sum_{j=0}^{k-1}(e^t)^j
=\boxed{\frac{e^t(1-e^{kt})}{k(1-e^t)}}.
\end{aligned}
\]
At $t=0$ the displayed quotient has the form $0/0$, but the defining sum gives $\boxed{M_X(0)=1}$; this also equals the limit of the quotient as $t\to0$.
\end{solution}

\begin{exercise}{7.18}
Let $X\sim\Geom(p)$, with $0<p<1$ and $q=1-p$, so that $\Prob(X=x)=pq^{x-1}$ for $x=1, 2, \ldots$. Show that
\[
M_X(t)=\frac{pe^t}{1-qe^t}, \qquad t<-\ln q,
\]
and use the MGF to find the mean and variance of $X$.
\end{exercise}
\begin{solution}
For a convergent geometric series we require $qe^t<1$, equivalently $t<-\ln q$. Hence
\[
M_X(t)=\sum_{x=1}^\infty e^{tx}pq^{x-1}
=pe^t\sum_{j=0}^\infty(qe^t)^j
=\boxed{\frac{pe^t}{1-qe^t}}.
\]
Differentiating in the interval of convergence gives
\[
M'_X(t)=\frac{pe^t}{(1-qe^t)^2}, \qquad
M''_X(t)=\frac{pe^t(1+qe^t)}{(1-qe^t)^3}.
\]
Since $1-q=p$,
\[
\Exp[X]=M'_X(0)=\boxed{\frac1p}, \qquad
\Var(X)=M''_X(0)-M'_X(0)^2
=\frac{1+q}{p^2}-\frac1{p^2}=\boxed{\frac{q}{p^2}}.
\]
\end{solution}

\begin{exercise}{7.19}
Let $X\sim\Poi(\mu)$, with $\mu>0$, so that $\Prob(X=x)=e^{-\mu}\mu^x/x!$ for $x=0, 1, 2, \ldots$. Show that $M_X(t)=e^{\mu(e^t-1)}$, and use the MGF to find the mean and variance of $X$.
\end{exercise}
\begin{solution}
Using the exponential series,
\[
M_X(t)=\sum_{x=0}^\infty e^{tx}\frac{e^{-\mu}\mu^x}{x!}
=e^{-\mu}\sum_{x=0}^\infty\frac{(\mu e^t)^x}{x!}
=\boxed{e^{\mu(e^t-1)}}.
\]
The MGF is finite for every $t\in\R$, and
\[
M'_X(t)=\mu e^tM_X(t), \qquad
M''_X(t)=(\mu e^t+\mu^2e^{2t})M_X(t).
\]
As $M_X(0)=1$,
\[
\Exp[X]=M'_X(0)=\boxed{\mu}, \qquad
\Var(X)=M''_X(0)-M'_X(0)^2
=\mu+\mu^2-\mu^2=\boxed{\mu}.
\]
\end{solution}

\begin{exercise}{7.20}
The MGF of a certain Poisson random variable $X$ is $M_X(t)=e^{9(e^t-1)}$. Find $\Prob(\mu-\sigma\le X\le\mu+\sigma)$, where $\mu=\Exp[X]$ and $\sigma=\sqrt{\Var(X)}$.
\end{exercise}
\begin{solution}
Comparing with the Poisson MGF identifies $X\sim\Poi(9)$, hence $\mu=9$ and $\sigma=3$. Since $X$ is integer-valued,
\[
\begin{aligned}
\Prob(\mu-\sigma\le X\le\mu+\sigma)
&=\Prob(6\le X\le12)\\
&=e^{-9}\sum_{x=6}^{12}\frac{9^x}{x!}
\approx\boxed{0.7601}.
\end{aligned}
\]
\end{solution}

\begin{exercise}{7.21}
Show that the MGF of a chi-squared random variable $X\sim\chi^2_\nu$ with $\nu$ degrees of freedom is
\[
M_X(t)=(1-2t)^{-\nu/2}.
\]
\end{exercise}
\begin{solution}
Its density is $f_X(x)=x^{\nu/2-1}e^{-x/2}/(2^{\nu/2}\Gamma(\nu/2))$ for $x>0$. Thus, for $t<1/2$,
\[
\begin{aligned}
M_X(t)
&=\frac1{2^{\nu/2}\Gamma(\nu/2)}
\int_0^\infty x^{\nu/2-1}e^{-(1/2-t)x}\,\dd x\\
&=\frac{\Gamma(\nu/2)}{2^{\nu/2}\Gamma(\nu/2)(1/2-t)^{\nu/2}}
=\boxed{(1-2t)^{-\nu/2}}.
\end{aligned}
\]
The gamma integral is finite precisely when $1/2-t>0$; this is the domain of the MGF.
\end{solution}

\begin{exercise}{7.22}
The MGF of a chi-squared random variable $X$ with $\nu$ degrees of freedom is $M_X(t)=(1-2t)^{-\nu/2}$ for $t<1/2$. Use it to show that $\Exp[X]=\nu$ and $\Var(X)=2\nu$.
\end{exercise}
\begin{solution}
Differentiation gives
\[
M'_X(t)=\nu(1-2t)^{-\nu/2-1}, \qquad
M''_X(t)=\nu(\nu+2)(1-2t)^{-\nu/2-2}.
\]
Therefore
\[
\Exp[X]=M'_X(0)=\boxed{\nu}, \qquad
\Var(X)=M''_X(0)-M'_X(0)^2
=\nu(\nu+2)-\nu^2=\boxed{2\nu}.
\]
\end{solution}

\begin{exercise}{7.24}
Let $X$ be a continuous random variable with density $f_X$. Suppose its MGF is finite on an open interval containing $0$. Recall that
\[
M_X(t)=\int_{-\infty}^{\infty}e^{tx}f_X(x)\,\dd x.
\]
Expand $e^{tx}$ in a Maclaurin series and integrate term by term to show that
\[
M_X(t)=1+\mu t+\mu'_2\frac{t^2}{2!}+\cdots+\mu'_r\frac{t^r}{r!}+\cdots,
\]
where $\mu'_r=\Exp[X^r]$ and $\mu=\Exp[X]$.
\end{exercise}
\begin{solution}
For $t$ sufficiently close to $0$, the MGF hypothesis permits term-by-term integration of the exponential series:
\[
\begin{aligned}
M_X(t)
&=\int_{-\infty}^{\infty}\left(\sum_{r=0}^\infty\frac{(tx)^r}{r!}\right)f_X(x)\,\dd x\\
&=\sum_{r=0}^\infty\frac{t^r}{r!}
\int_{-\infty}^{\infty}x^rf_X(x)\,\dd x
=\sum_{r=0}^\infty\frac{\mu'_r}{r!}t^r.
\end{aligned}
\]
Since
\[
\mu'_0=\int_{-\infty}^{\infty}f_X(x)\,\dd x=1,\qquad\mu'_1=\mu,
\]
this is the required expansion. Equivalently, $M_X^{(r)}(0)=\mu'_r$.
\end{solution}

\Needspace{14\baselineskip}
\begin{exercise}{8.1}
Define a suitable population for each of the following samples.
\begin{exercisesubparts}
\item Customers contacted by a car manufacturer for feedback after their vehicles were serviced at its authorized service center.
\item Ten randomly selected college students, of whom four are girls.
\item Total marks collected from 100 graduate students at a university.
\item Fifty observations made to determine the lifetime of a particular brand of truck tire.
\end{exercisesubparts}
\end{exercise}
\begin{solution}
\begin{exercisesubparts}
\item The population is all customers whose vehicles are serviced at the manufacturer's authorized service center (in the period of interest).
\item The population is all students at the college from which the ten were randomly selected, classified by gender.
\item The population is the set of graduate students at that university, with each student's total marks as the measured characteristic.
\item The population is the lifetimes of all tires of that brand under the specified usage conditions.
\end{exercisesubparts}
\end{solution}

\begin{exercise}{8.2}
The waiting times, in minutes, of ten patients at a doctor's office are
\[
5, \ 11, \ 9, \ 5, \ 10, \ 15, \ 6, \ 10, \ 5, \ 10.
\]
Treat these as a random sample and find
\begin{exercisesubparts}
\item the mean.
\item the median.
\item the mode.
\end{exercisesubparts}
\end{exercise}
\begin{solution}
The ordered data are $5, 5, 5, 6, 9, 10, 10, 10, 11, 15$, with $n=10$ and total $86$.
\begin{exercisesubparts}
\item $\displaystyle\bar x=\frac{86}{10}=\boxed{8.6\text{ min}}$.
\item The middle values are $9$ and $10$, so $\displaystyle\tilde x=\frac{9+10}{2}=\boxed{9.5\text{ min}}$.
\item Both $5$ and $10$ occur three times, the greatest frequency. The modes are $\boxed{5\text{ and }10\text{ min}}$.
\end{exercisesubparts}
\end{solution}

\begin{exercise}{8.3}
Reaction times, in seconds, for nine subjects given a stimulant were recorded as
\[
2.5, \ 3.6, \ 3.1, \ 4.3, \ 2.9, \ 2.3, \ 2.6, \ 4.1, \ 3.4.
\]
Calculate
\begin{exercisesubparts}
\item the mean.
\item the median.
\end{exercisesubparts}
\end{exercise}
\begin{solution}
\begin{exercisesubparts}
\item Since $\displaystyle \sum x_i=28.8$, $\displaystyle\bar x=\frac{28.8}{9}=\boxed{3.2\text{ s}}$.
\item The ordered values are $2.3, 2.5, 2.6, 2.9, 3.1, 3.4, 3.6, 4.1, 4.3$. The fifth is the middle observation, so $\boxed{\tilde x=3.1\text{ s}}$.
\end{exercisesubparts}
\end{solution}

\Needspace{14\baselineskip}
\begin{exercise}{8.5}
The numbers of incorrect answers on a true-false competency test for a random sample of 15 students are
\[
2, \ 1, \ 3, \ 0, \ 1, \ 3, \ 6, \ 0, \ 3, \ 3, \ 5, \ 2, \ 1, \ 4, \ 2.
\]
Find
\begin{exercisesubparts}
\item the mean.
\item the median.
\item the mode.
\end{exercisesubparts}
\end{exercise}
\begin{solution}
The sum is $36$. In ascending order the data are $0, 0, 1, 1, 1, 2, 2, 2, 3, 3, 3, 3, 4, 5, 6$.
\begin{exercisesubparts}
\item $\displaystyle\bar x=\frac{36}{15}=\boxed{2.4}$.
\item The eighth of $15$ ordered values is $2$, so $\boxed{\tilde x=2}$.
\item The value $3$ occurs four times, more than any other value. Thus the mode is $\boxed{3}$.
\end{exercisesubparts}
\end{solution}

\begin{exercise}{8.6}
The numbers of medical leaves claimed in a year by ten employees are
\[
5, \ 12, \ 3, \ 7, \ 11, \ 1, \ 85, \ 5, \ 9, \ 5.
\]
Find the mean, median, and mode. Which gives the best measure of the center? Explain.
\end{exercise}
\begin{solution}
The observations sum to $143$, and their ordered list is $1, 3, 5, 5, 5, 7, 9, 11, 12, 85$. Therefore
\[
\bar x=\boxed{14.3}, \qquad
\tilde x=\frac{5+7}{2}=\boxed{6}, \qquad
\text{mode}=\boxed{5}.
\]
The unusually large observation $85$ raises the mean far above the bulk of the data. The \emph{median} $6$ is the preferable measure of center here because it is resistant to this extreme value.
\end{solution}

\begin{exercise}{8.7}
The examination scores of twenty students selected from a city school are
\[
\begin{gathered}
45, \ 57, \ 68, \ 34, \ 50, \ 32, \ 89, \ 47, \ 97, \ 67, \\
79, \ 84, \ 43, \ 35, \ 68, \ 55, \ 72, \ 63, \ 68, \ 49.
\end{gathered}
\]
Calculate
\begin{exercisesubparts}
\item the mean.
\item the mode.
\end{exercisesubparts}
\end{exercise}
\begin{solution}
\begin{exercisesubparts}
\item The score total is $1202$, so $\displaystyle\bar x=\frac{1202}{20}=\boxed{60.1}$.
\item The score $68$ appears three times; all other scores appear once. The mode is $\boxed{68}$.
\end{exercisesubparts}
\end{solution}

\Needspace{17\baselineskip}
\begin{exercise}{8.8}
Phosphates in household detergents can contribute to water pollution. The following are phosphate amounts (grams per prescribed load) for twelve sampled detergents:
\[
\begin{array}{lc@{\qquad\qquad}lc}
\toprule
\text{Detergent}&\text{Grams}&\text{Detergent}&\text{Grams}\\
\midrule
\text{A \& P Blue Sail}&48&\text{Ajax}&31\\
\text{Dash}&47&\text{Sears}&30\\
\text{Concentrated All}&42&\text{Fab}&29\\
\text{Cold Water All}&42&\text{Cold Power}&29\\
\text{Breeze}&41&\text{Bold}&29\\
\text{Oxydol}&34&\text{Rinso}&26\\
\bottomrule
\end{array}
\]
For the phosphate data, find
\begin{exercisesubparts}
\item the mean.
\item the median.
\item the mode.
\end{exercisesubparts}
\end{exercise}
\begin{solution}
There are $12$ observations, with total $428$. In ascending order the values are
\[
26, \ 29, \ 29, \ 29, \ 30, \ 31, \ 34, \ 41, \ 42, \ 42, \ 47, \ 48.
\]
\begin{exercisesubparts}
\item $\displaystyle\bar x=\frac{428}{12}=\boxed{\frac{107}{3}\approx35.67\text{ g}}$.
\item The middle values are $31$ and $34$, so $\displaystyle\tilde x=\frac{31+34}{2}=\boxed{32.5\text{ g}}$.
\item The most frequent amount is $29$ (three occurrences). The mode is $\boxed{29\text{ g}}$.
\end{exercisesubparts}
\end{solution}

\begin{exercise}{8.9}
The waiting times, in minutes, of ten patients are $5, 11, 9, 5, 10, 15, 6, 10, 5, 10$. Treat them as a random sample and find
\begin{exercisesubparts}
\item the range.
\item the sample standard deviation.
\end{exercisesubparts}
\end{exercise}
\begin{solution}
\noindent\textup{(a)} The smallest and largest times are $5$ and $15$, so $\boxed{r=15-5=10\text{ min}}$.

\noindent\textup{(b)} For $n=10$, the sum of observations is $86$ and the sum of squares is $838$. The computational formula gives
\[
s^2=\frac{838-86^2/10}{10-1}=\frac{164}{15}\approx10.9333\text{ min}^2,
\qquad \boxed{s=\sqrt{\frac{164}{15}}\approx3.307\text{ min}}.
\]
\end{solution}

\Needspace{14\baselineskip}
\begin{exercise}{8.10}
For the reaction times (seconds)
\[
2.5, \ 3.6, \ 3.1, \ 4.3, \ 2.9, \ 2.3, \ 2.6, \ 4.1, \ 3.4,
\]
calculate
\begin{exercisesubparts}
\item the range.
\item the sample variance using the definition $\displaystyle s^2=\frac{1}{n-1}\sum_{i=1}^{n}(x_i-\bar x)^2$.
\end{exercisesubparts}
\end{exercise}
\begin{solution}
\noindent\textup{(a)} $\boxed{r=4.3-2.3=2.0\text{ s}}$.


\noindent\textup{(b)} The sample mean is $\bar x=3.2$. By the definition,
\[
s^2=\frac18\sum_{i=1}^9(x_i-\bar x)^2
=\frac18\sum_{i=1}^9(x_i-3.2)^2
=\frac{3.98}{8}=\boxed{0.4975\text{ s}^2}.
\]
\end{solution}

\Needspace{14\baselineskip}
\begin{exercise}{8.11}
For the sample of numbers of incorrect test answers
\[
2, \ 1, \ 3, \ 0, \ 1, \ 3, \ 6, \ 0, \ 3, \ 3, \ 5, \ 2, \ 1, \ 4, \ 2,
\]
calculate the sample variance using
\begin{exercisesubparts}
\item the definition $\displaystyle s^2=\frac{1}{n-1}\sum_{i=1}^{n}(x_i-\bar x)^2$.
\item the computational formula of Theorem 8.1, $\displaystyle s^2=\frac{n\sum_{i=1}^{n}x_i^2-(\sum_{i=1}^{n}x_i)^2}{n(n-1)}$.
\end{exercisesubparts}
\end{exercise}
\begin{solution}
Here $n=15$, $\bar x=2.4$, and the data give
\[
\sum_{i=1}^{15}x_i=36, \qquad \sum_{i=1}^{15}x_i^2=128.
\]
\begin{exercisesubparts}
\item By the definition, group the equal observations:
\[
\begin{aligned}
\sum_{i=1}^{15}(x_i-\bar x)^2
&=2(0-2.4)^2+3(1-2.4)^2+3(2-2.4)^2\\
&\quad+4(3-2.4)^2+(4-2.4)^2+(5-2.4)^2+(6-2.4)^2\\
&=41.6.
\end{aligned}
\]
Therefore
\[
s^2=\frac{41.6}{14}=\boxed{\frac{104}{35}\approx2.9714}.
\]
\item By Theorem 8.1,
\[
\begin{aligned}
s^2&=\frac{n\sum_{i=1}^n x_i^2-(\sum_{i=1}^n x_i)^2}{n(n-1)}\\
&=\frac{15(128)-36^2}{15(14)}
=\boxed{\frac{104}{35}\approx2.9714}.
\end{aligned}
\]
\end{exercisesubparts}
\end{solution}

\Needspace{11\baselineskip}
\begin{exercise}{8.12}
The saturated-fat contents, in grams, of eight randomly selected cereal bars are
\[
0.65, \ 0.72, \ 0.45, \ 0.55, \ 0.58, \ 0.39, \ 0.68, \ 0.52.
\]
Calculate
\begin{exercisesubparts}
\item the sample mean.
\item the sample variance.
\end{exercisesubparts}
\end{exercise}
\begin{solution}
For $n=8$, the sum of observations is $4.54$ and the sum of squares is $2.6672$.
\begin{exercisesubparts}
\item $\displaystyle\bar x=\frac{4.54}{8}=\boxed{0.5675\text{ g}}$.
\item The computational formula yields
\[
s^2=\frac{2.6672-(4.54)^2/8}{7}
=\boxed{0.0129643\text{ g}^2}.
\]
\end{exercisesubparts}
\end{solution}

\Needspace{12\baselineskip}
\begin{exercise}{8.13}
The heights, in meters, of twenty randomly selected college seniors are
\[
\begin{array}{ccccc}
\toprule
1.53&1.48&1.64&1.70&1.56\\
1.79&1.69&1.74&1.79&1.73\\
1.43&1.68&1.64&1.75&1.76\\
1.59&1.68&1.74&1.73&1.71\\
\bottomrule
\end{array}
\]
Calculate the sample standard deviation.
\end{exercise}
\begin{solution}
For $n=20$, the data give $\displaystyle \sum x_i=33.36$ and $\displaystyle \sum x_i^2=55.8414$, so $\bar x=1.668$ m. The computational formula gives
\[
\begin{aligned}
s^2&=\frac{55.8414-(33.36)^2/20}{19}
\approx0.01036421\text{ m}^2, \\
s&=\boxed{\sqrt{s^2}\approx0.1018\text{ m}}.
\end{aligned}
\]
\end{solution}

\begin{exercise}{8.14}
For a sample $x_1, \ldots, x_n$ with $n\ge2$, prove that
\begin{exercisesubparts}
\item adding or subtracting a constant $c$ from every observation leaves the sample variance unchanged.
\item multiplying every observation by a constant $c$ multiplies the sample variance by $c^2$.
\end{exercisesubparts}
\end{exercise}
\begin{solution}
Write
\[
s_x^2=\frac{1}{n-1}\sum_{i=1}^n(x_i-\bar x)^2.
\]
\begin{exercisesubparts}
\item If $y_i=x_i+c$, then $\bar y=\bar x+c$, so $y_i-\bar y=x_i-\bar x$. Thus $\boxed{s_y^2=s_x^2}$. Replacing $c$ by $-c$ proves the subtraction case.
\item If $y_i=cx_i$, then $\bar y=c\bar x$ and $y_i-\bar y=c(x_i-\bar x)$. Consequently,
\[
s_y^2=\frac1{n-1}\sum_{i=1}^n c^2(x_i-\bar x)^2
=\boxed{c^2s_x^2}.
\]
\end{exercisesubparts}
\end{solution}

\Needspace{10\baselineskip}
\begin{exercise}{8.15}
First calculate the sample variance of $8, 10, 16, 18, 24, 25$. Then use the variance rules for translation and scaling to find the variances of
\begin{exercisesubparts}
\item $16, 20, 32, 36, 48, 50$.
\item $12, 14, 20, 22, 28, 29$.
\end{exercisesubparts}
\end{exercise}
\begin{solution}
For the original sample, $n=6$, $\displaystyle \sum x_i=101$, and $\displaystyle \sum x_i^2=1945$. Hence
\[
s_x^2=\frac{1945-101^2/6}{5}=\boxed{\frac{1469}{30}\approx48.9667}.
\]
\begin{exercisesubparts}
\item Each value is doubled, so $\boxed{s_{2x}^2=4s_x^2=\frac{2938}{15}\approx195.8667}$.
\item Each value is increased by $4$, so $\boxed{s_{x+4}^2=s_x^2=\frac{1469}{30}\approx48.9667}$.
\end{exercisesubparts}
\end{solution}

\begin{exercise}{8.16}
In twelve one-day cricket matches during the 2014--15 season, a university team captain scored
\[
84, \ 25, \ 74, \ 53, \ 40, \ 31, \ 64, \ 71, \ 18, \ 63, \ 88, \ 49.
\]
Find
\begin{exercisesubparts}
\item the mean score.
\item the median score.
\end{exercisesubparts}
\end{exercise}
\begin{solution}
\begin{exercisesubparts}
\item The total is $660$, so $\displaystyle\bar x=\frac{660}{12}=\boxed{55}$ runs.
\item The ordered scores are $18, 25, 31, 40, 49, 53, 63, 64, 71, 74, 84, 88$. The sixth and seventh are $53$ and $63$, so $\displaystyle\tilde x=\frac{53+63}{2}=\boxed{58}$ runs.
\end{exercisesubparts}
\end{solution}

\begin{exercise}{8.17}
All possible samples of size $25$ are drawn from a normal population with mean $60$ and standard deviation $9$. Find the probability that the sample mean $\bar X$ lies between $\mu_{\bar X}-0.96\sigma_{\bar X}$ and $\mu_{\bar X}+0.25\sigma_{\bar X}$. Assume the sample means can be measured to any degree of accuracy.
\end{exercise}
\begin{solution}
Since the population is normal, $\bar X\sim\Normal(60, 9^2/25)$ exactly; thus $\sigma_{\bar X}=9/5$. Standardization gives
\[
\begin{aligned}
&\Prob(\mu_{\bar X}-0.96\sigma_{\bar X}<\bar X<\mu_{\bar X}+0.25\sigma_{\bar X})\\
&\qquad=\Prob(-0.96<Z<0.25)
=\Phi(0.25)-\Phi(-0.96)
\approx\boxed{0.4302}.
\end{aligned}
\]
\end{solution}

\Needspace{8\baselineskip}
\begin{exercise}{8.18}
For random samples of size $49$ from a large or infinite population, the standard deviation of the sampling distribution of the mean is $3$. How large must the sample size be to reduce this standard deviation to $1.5$?
\end{exercise}
\begin{solution}
The standard error is $\sigma_{\bar X}=\sigma/\sqrt n$. With the same population variance,
\[
\frac{1.5}{3}=\sqrt{\frac{49}{n}}
\iff n=49\left(\frac3{1.5}\right)^2
=\boxed{196}.
\]
Thus quadrupling the sample size halves the standard error.
\end{solution}

\begin{exercise}{8.19}
An automobile battery has mean lifetime $1500$ days and standard deviation $250$ days. How does the variance of the sample mean change when the sample size is
\begin{exercisesubparts}
\item increased from $49$ to $225$?
\item decreased from $625$ to $36$?
\end{exercisesubparts}
\end{exercise}
\begin{solution}
Independence of sample observations gives $\Var(\bar X)=\sigma^2/n=62500/n$ (days$^2$).
\begin{exercisesubparts}
\item $\displaystyle\Var(\bar X_{49})=\frac{62500}{49}\approx1275.5102$ and $\displaystyle\Var(\bar X_{225})=\frac{62500}{225}\approx\boxed{277.7778}$. The variance decreases by the factor $49/225$.
\item $\displaystyle\Var(\bar X_{625})=\frac{62500}{625}=100$ and $\displaystyle\Var(\bar X_{36})=\frac{62500}{36}\approx\boxed{1736.1111}$. The variance increases by the factor $625/36$.
\end{exercisesubparts}
\end{solution}

\Needspace{14\baselineskip}
\begin{exercise}{8.21}
A soft-drink machine dispenses an average of $240$ mL per drink, with standard deviation $15$ mL. From time to time, a random sample of $40$ drinks is taken. The machine is considered satisfactory if the sample mean lies within $\mu_{\bar X}\pm2\sigma_{\bar X}$; otherwise it is adjusted. An observed sample mean is $236$ mL. Was the decision not to adjust the machine reasonable?
\end{exercise}
\begin{solution}
For a random sample of size $40$,
\[
\mu_{\bar X}=240, \qquad
\sigma_{\bar X}=\frac{15}{\sqrt{40}}\approx2.3717\text{ mL}.
\]
The acceptance interval is
\[
\left[240-\frac{30}{\sqrt{40}}, \ 240+\frac{30}{\sqrt{40}}\right]
\approx[235.257, 244.743]\text{ mL}.
\]
The observed mean $236$ mL lies within the acceptance interval. Therefore no adjustment is required under the specified rule; the sample mean need not equal $240$ exactly.
\end{solution}

\begin{exercise}{8.25}
The time taken to complete a project at a real estate company has mean $18$ months and standard deviation $3$ months; assume project completion times are approximately normal. Find
\begin{exercisesubparts}
\item the probability that the mean completion time for four independent projects lies between $16$ and $19$ months.
\item the value of $\bar x$ above which $20\%$ of sample means of size five would fall.
\end{exercisesubparts}
\end{exercise}
\begin{solution}
\noindent\textup{(a)} For $n=4$, $\bar X$ is approximately $\Normal(18, 9/4)$, with standard error $3/\sqrt4=1.5$. Thus
\[
\begin{aligned}
\Prob(16<\bar X<19)
&\approx\Phi\left(\frac{19-18}{1.5}\right)-\Phi\left(\frac{16-18}{1.5}\right)\\
&=\Phi\left(\frac23\right)-\Phi\left(-\frac43\right)
\approx\boxed{0.6563}.
\end{aligned}
\]
\noindent\textup{(b)} For $n=5$, the upper-tail probability $0.20$ corresponds to the 80th percentile, $z_{0.80}=\Phi^{-1}(0.80)\approx0.8416$. Therefore
\[
\boxed{\bar x=18+\frac3{\sqrt5}z_{0.80}\approx19.13\text{ months}}.
\]
\end{solution}

\begin{exercise}{8.35}
An electrical firm makes light bulbs whose lifetimes are approximately normal with mean $800$ hours and standard deviation $40$ hours. A random sample of $16$ bulbs has mean lifetime $\bar x=775$ hours. Does this observation cast doubt on the premise $\mu=800$? Calculate $\Prob(\bar X\le775)$ first when $\mu=800$, and then when $\mu=760$, keeping the standard deviation $40$ hours.
\end{exercise}
\begin{solution}
With $n=16$, the sampling standard deviation is $40/\sqrt{16}=10$ hours. Under the assumption $\mu=800$,
\[
\Prob(\bar X\le775)\approx\Phi\left(\frac{775-800}{10}\right)
=\Phi(-2.5)\approx\boxed{0.0062}.
\]
Under the alternative assumption $\mu=760$, with the same standard deviation,
\[
\Prob(\bar X\le775)\approx\Phi\left(\frac{775-760}{10}\right)
=\Phi(1.5)\approx\boxed{0.9332}.
\]
A sample mean this low is quite unusual when $\mu=800$, but is not unusual when $\mu=760$. The observation therefore casts doubt on the original mean of $800$ hours; it does not by itself determine the true population mean.
\end{solution}

\begin{exercise}{8.60}
Five retail shops are selected at random from a city, and their quoted prices for a home appliance are $\$305$, $\$312$, $\$296$, $\$304$, and $\$307$. Find the sample variance of these quotations.
\end{exercise}
\begin{solution}
The sample size is $5$, and
\[
\sum_{i=1}^5x_i=1524, \qquad
\sum_{i=1}^5x_i^2=464650, \qquad
\bar x=\frac{1524}{5}=304.8.
\]
Therefore
\[
s^2=\frac{464650-1524^2/5}{5-1}
=\boxed{33.7\text{ dollars}^2}.
\]
The denominator is $n-1=4$, since the quotations constitute a sample.
\end{solution}

\begin{exercise}{8.72}
A random variable $X$ is normally distributed with mean $20$ and variance $9$. Determine the smallest positive integer sample size $n$ for which the mean $\bar X$ of an independent random sample satisfies
\[
\Prob(19.9\le\bar X\le20.1)\ge0.95.
\]
\end{exercise}
\begin{solution}
Because $X$ is normal, $\bar X\sim\Normal(20, 9/n)$ exactly. By symmetry and standardization,
\[
\Prob(19.9\le\bar X\le20.1)
=\Prob\left(-\frac{0.1\sqrt n}{3}\le Z\le\frac{0.1\sqrt n}{3}\right)
=2\Phi\left(\frac{0.1\sqrt n}{3}\right)-1.
\]
The required probability is at least $0.95$ precisely when
\[
\frac{0.1\sqrt n}{3}\ge\Phi^{-1}(0.975)\approx1.959964,
\quad\text{that is,}\quad
n\ge\left(\frac{3\Phi^{-1}(0.975)}{0.1}\right)^2
\approx3457.313.
\]
Thus the minimum integer sample size is $\boxed{n=3458}$.
\end{solution}

\end{document}
